\newpage
\section{Introduction}
\label{sec:introduction}
\subsection{Contexte du projet}
% Présentez ici le contexte général du projet
Dans le cadre de notre cours de Génie Logiciel du Projet en Licence 2 Informatique à CY Cergy-Paris Université, on nous a demandé de concevoir et développer un jeu en Java. Notre groupe a choisi de réaliser une version modifiée du jeu d'échecs chinois, qu'on a appelé \textbf{CHESS}. Ce choix vient de notre intérêt commun pour les jeux de stratégie, et ça nous permettait de travailler sur quelque chose de concret tout en appliquant ce qu'on a vu en cours, comme la séparation entre le moteur de jeu et l'interface graphique.
\paragraph{} Les éléments présentés dans ce rapport correspondent à l'ensemble du travail réalisé par notre groupe dans le cadre du projet Génie Logiciel.
\subsection{Objectif du projet}
% Décrivez l'objectif principal du logiciel
L'objectif était de faire un jeu jouable avec une interface graphique, un moteur qui gère toutes les règles correctement, et un bot pour jouer contre l'ordinateur. On voulait que le jeu soit complet et fonctionnel, pas juste une démonstration.
\subsection{Organisation du rapport}
% Expliquez comment le rapport est structuré
Dans ce rapport, on présente les sections suivantes :
\begin{enumerate}
\item Les spécifications du projet.
\item La conception et le développement du logiciel.
\item Le manuel utilisateur.
\item Le déroulement du projet en équipe.
\item La conclusion et les perspectives d'amélioration.
\end{enumerate}
